\documentclass[10pt]{article}
\usepackage[utf8]{inputenc}
\usepackage[T1]{fontenc}
\usepackage{charter} 
\usepackage[margin=0.9in]{geometry}
\usepackage{hyperref}
\usepackage{enumitem}

\pagestyle{empty}

\begin{document}

\begin{raggedright}

\textbf{Olga C. Santos} \\
PhyUM Research Center, Department of Artificial Intelligence, UNED \\
ocsantos@dia.uned.es \\

\vspace{0.2cm}

\textbf{Concepción Labra} \\
Department of Artificial Intelligence, Computer Science School, UNED \\
clabra@dia.uned.es \\

\end{raggedright}

\vspace{0.4cm}

\noindent
To the Editor,

\vspace{0.4cm}

\noindent

We are submitting the paper entitled \textbf{``Theory-Based Interpretability in Deep Knowledge Tracing via Grounded Transformers''} for consideration for publication in the Special Issue \textbf{``Artificial Intelligence and Educational Data Mining for Personalised Learning and Adaptive Instruction''} of the journal \textit{Applied Sciences}.

\subsection*{Significance and Context}
Knowledge Tracing is a foundational task in Educational Data Mining and a cornerstone of Intelligent Tutoring Systems. While deep learning models have established state-of-the-art predictive performance in this domain, their ``black-box'' nature remains a critical barrier to adoption in real-world educational settings where interpretability and trust are essential. 

Our work introduces \textbf{gTransformer}, a novel grounded Transformer model that bridges this gap through \textit{Representational Grounding}. By anchoring high-capacity deep learning models to reference models such as Bayesian Knowledge Tracing based on pedagogically sound principles, we demonstrate that competitive accuracy can coexist with intrinsic interpretability.

\subsection*{Key Contributions and Findings}
\begin{itemize}[leftmargin=*, label=$\bullet$, itemsep=0pt, topsep=2pt]
    \item \textbf{Triple Validation}: We propose a rigorous validation approach encompassing structural alignment, semantic consistency, and functional confidence to prove that the model genuinely internalizes pedagogical concepts rather than merely making predictions.
    \item \textbf{Performance vs. Interpretability}: Our results show that gTransformer achieves results competitive with state-of-the-art architectures, outperforming population-level theory-based models such as BKT by a significant \textbf{19.9\% AUC gain}, while incurring a minimal performance cost (3.9\%) relative to black-box configurations.
    \item \textbf{Context-Aware Personalization}: We demonstrate that gTransformer captures longitudinal learning dynamics that traditional Markovian models overlook. This enables differentiated diagnostics for students with identical response scores but different histories, providing a foundation for context-aware personalization.
\end{itemize}

\subsection*{Fit with Journal and Special Issue Scope}
This manuscript directly addresses the Special Issue's core themes. By leveraging advanced Deep Learning to improve education through ``early identification of student needs'' and ``data-driven personalization,'' our work contributes a practical, theory-grounded tool for the GED and Learning Analytics communities.

\subsection*{Mandatory Statements}
We confirm that neither the manuscript nor any parts of its content are currently under consideration for publication with or published in another journal. All authors have approved the manuscript and agree with its submission to \textit{Applied Sciences}.

\vspace{0.4cm}

\noindent
Sincerely,

\vspace{0.4cm}

\noindent
Olga C. Santos \\
Concepción Labra 

\end{document}
